\documentclass[12pt]{article}

\usepackage{amsfonts} % To do stylized R and stuff
\usepackage{amsmath}  % For the question* enviroment

% Margin - 1 inch on all sides
\usepackage[letterpaper]{geometry}
\usepackage{times}
\geometry{top=1.0in, bottom=1.0in, left=1.0in, right=1.0in}

%
% Doublespacing
%
% \usepackage{setspace}
% \doublespacing        % I am not going to use this here but it is good to
                        % know about

%
% Rotating tables (e.g. sideways when too long)
%
\usepackage{rotating}

%
% Fancy-header package to modify header/page numbering (insert last name)
%
\usepackage{fancyhdr}
\pagestyle{fancy}
\lhead{}            % Makes the left header empty
\chead{}            % Makes the center header empty
\rhead{\thepage}    % Puts a page number of the right header
\lfoot{}            % Makes left footer empty
\cfoot{}            % Makes center footer empty
\rfoot{}            % Makes right footer empty
\renewcommand{\headrulewidth}{0pt}  % Removes the horizontal line under the 
                                    % header.
\renewcommand{\footrulewidth}{0pt}  % Removes the horizontal line above the 
                                    % footer.
%To make sure we actually have header 0.5in away from top edge
%12pt is one-sixth of an inch. Subtract this from 0.5in to get headsep value
\setlength\headsep{0.333in}

%
% Makes it so that section numbers are bold, but not the text
%
\usepackage{titlesec}
\titleformat{\section}
  {\large}                  % Smaller than default \Large
  {\bfseries\thesection}    % Bold number
  {1em}
  {}

%%%%%%%%%%%%%%%%%
% Begin document
%%%%%%%%%%%%%%%%%
\begin{document}
\begin{flushleft}

%Title
\begin{center}
    \bfseries\LARGE Linear Algebra Exam 3 Study Guide
\end{center}

%Changes paragraph indentation to 0.5in
\setlength{\parindent}{0.5in}
%Begin body of paper here

\section{\textbf{Section 5.2}: Inner Product Spaces}
\begin{itemize}
    \item An \textbf{inner product} (denoted as $\langle \vec{u},\vec{v} 
            \rangle$) is literally just any operation that satifies the 
            following axioms (let $\vec{u}$, $\vec{v}$, \& $\vec{w}$ be in a 
            vector space $V$, and $c \in \mathbb{R}$):
            \begin{enumerate}
                \item $\langle \vec{u},\vec{v} \rangle = \langle \vec{v},
                    \vec{u} \rangle$
                \item $\langle \vec{u},\vec{v} + \vec{w} \rangle = \langle 
                    \vec{u}, \vec{v} \rangle + \langle \vec{u},\vec{w} 
                    \rangle$
                \item $c \langle \vec{u},\vec{v} \rangle = \langle c\vec{u},
                    \vec{v} \rangle = \langle \vec{u},c\vec{v} \rangle$
                \item $\langle \vec{u},\vec{u} \rangle > 0$ \& $\langle 
                    \vec{u}, \vec{v} \rangle = 0 $ iff $\vec{u}=0$
            \end{enumerate}
            Note: The dot product is an example of an inner product (it is also 
            called the Euclidean inner product)
    \item Note: An inner product can also be done on matrices, and even
            definite integrals
    \item \textbf{Inner product space}: a vector space with an inner product
            defined
    \item \textbf{Length of a vector}: $||\vec{u}|| = \sqrt{\langle \vec{u},
            \vec{v} \rangle}$
    \item \textbf{Distance between 2 vectors}: $d(\vec{u},\vec{v})=||\vec{u}-
            \vec{v}||$
    \item \textbf{The angle between 2 non-zero vectors}:
            \begin{align*}
                \cos(\theta) = \frac{\langle \vec{u},\vec{v} \rangle}
                {||\vec{u}||||\vec{v}||} \text{, \& } 0 \leq \theta \leq \pi
            \end{align*}
    \item \textbf{Orthogonal}: $\langle \vec{u}, \vec{v} \rangle = 0$
    \item \textbf{Unit vector eq}: $\vec{u} = \frac{\vec{v}}{||\vec{v}||}$ (or 
            to normalize)
    \item \textbf{Orthogonal Projections}: a vector broken down into 2 parts
            \begin{itemize}
                \item $\vec{\text{u}} = \hat{\text{u}} + \vec{z}$
                        \begin{itemize}
                            \item $\vec{\text{u}}$: a vector
                            \item $\hat{\text{u}}$: the projection of 
                                    $\vec{\text{u}}$ on another vector 
                                    $\vec{v}$. Think of it has the shadow that 
                                    $\vec{\text{u}}$ casts on all scalar 
                                    multiples of $\vec{v}$
                            \item $\vec{z}$: The vector orthogonal to $\vec{v}$
                        \end{itemize}
                \item $\hat{\text{u}} = \text{proj}_v \vec{\text{u}} = 
                        \frac{\langle \vec{\text{u}},\vec{v} \rangle}{\langle 
                        \vec{v}, \vec{v} \rangle}\vec{v}$
                \item $\vec{z} = \vec{\text{u}} - \hat{\text{u}}$
            \end{itemize}
\end{itemize}

\section{\textbf{Section 5.3}: Orthogonal Basis - Gram-Schmidt}
\begin{itemize}
    \item Vector spaces can have many basis, like how $\mathbb{R}^3$ has the
            basis $B = \{ (1,0,0), (0,1,0), (0,0,1) \}$ (the standard basis for
            $\mathbb{R}^3$). Standard basis have some important characteristics,
            like how all the vectors in it are mutually orthogonal
            \begin{align*}
                (1,0,0) \cdot (0,1,0) &= 0 \\
                (1,0,0) \cdot (0,0,1) &= 0 \\
                (0,1,0) \cdot (0,0,1) &= 0
            \end{align*}
    \item \textbf{Orthogonal set}: $S = \{ \vec{v}_1, \vec{v}_1, \ldots, 
            \vec{v}_n \} \in \mathbb{R}^n$ such that $\langle \vec{u}_i,
            \vec{u}_j \rangle = 0$ $\forall$ $i \neq j$. Basically, all vectors 
            are orthogonal to each other
    \item \textbf{Othonormal set}: $S = \{ \vec{v}_1, \vec{v}_1, \ldots, 
            \vec{v}_n \} \in \mathbb{R}^n$ such that: $\langle \vec{u}_i,
            \vec{u}_j \rangle = 0$ $\forall$ $i \neq j$, and every vector is a
            unit vector.
    \item \textbf{Gram-Schmidt Orthonormalization Algorithm}: Algorithm to find 
            orthonormal bases, where you start with a basis, you turn it into an
            orthonormal basis, then normalize each vector.
            \begin{enumerate}
                \item Let $B = \{ \vec{v}_1, \vec{v}_1, \ldots, \vec{v}_n \}$ be
                        a basis for an inner product space $V$
                \item Let $B' = \{ \vec{w}_1, \vec{w}_1, \ldots, \vec{w}_n \}$, 
                        where
                        \begin{align*}
                            \vec{w}_1 &= \vec{v}_1                              \\
                            \vec{w}_2 &= \vec{v}_2 - \frac{\langle \vec{v}_2,
                            \vec{w}_1 \rangle}{\langle \vec{w}_1, \vec{w}_1 
                            \rangle} \vec{w}_1                                  \\
                            \vec{w}_3 &= \vec{v}_3 - \frac{\langle \vec{v}_3,
                            \vec{w}_1 \rangle}{\langle \vec{w}_1, \vec{w}_1 
                            \rangle} \vec{w}_1 - \frac{\langle \vec{v}_3,
                            \vec{w}_1 \rangle}{\langle \vec{w}_1, \vec{w}_1 
                            \rangle} \vec{w}_1                                  \\
                            \vdots                                              \\
                            \vec{w}_n &= \vec{v}_n - \frac{\langle \vec{v}_n,
                            \vec{w}_1 \rangle}{\langle \vec{w}_1, \vec{w}_1 
                            \rangle} \vec{w}_1 - \frac{\langle \vec{v}_n, 
                            \vec{w}_1 \rangle} {\langle \vec{w}_1, \vec{w}_1 
                            \rangle} \vec{w}_1 - \dots - \frac{\langle 
                            \vec{v}_n, \vec{w}_{n-1} \rangle} {\langle 
                            \vec{w}_{n-1}, \vec{w}_{n-1} \rangle} \vec{w}_{n-1}
                        \end{align*}
                        Then $B'$ is an orthonormal basis for $V$
                \item Let $\vec{u}_i = \frac{\vec{w}_i}{||\vec{w}_i||}$. Then
                        $B'' = \{ \vec{u}_1, \vec{u}_1, \ldots, \vec{u}_n \}$ is
                        an orthonormal basis for $V$
            \end{enumerate}
            Also, $\text{span}\{ \vec{v}_1, \vec{v}_1, \ldots, \vec{v}_n \} =
            \text{span}\{ \vec{u}_1, \vec{u}_1, \ldots, \vec{u}_n \}$ for $k =
            1,2,\ldots, n$
\end{itemize}

\section{\textbf{Section 6.1}: Introductions to Linear Transformations}
\begin{itemize}
    \item Get back to this 
\end{itemize}

\end{flushleft}
\end{document}